\section{Gaussian marginal likelihood and profiling the difference mean}
\subsection{Likelihood selection for fixed $\mu$}
A probabilistic interpretation of the penalty assumes observation
noise $\varepsilon\sim N(0,\sigma^2I_n)$ and stochastic differences
$D_d\tau=\mu\mathbf1+\eta$ with
$\eta\sim N(0,\sigma^2\lambda^{-1}I_{n-d})$.
The $d$ unpenalized polynomial coordinates act as nuisance
fixed effects. In the $r=n-d$ positive-penalty singular coordinates
$z_j=v_j^\top y$, write $p^{(1)}_t=\binom td$,
$b_j=v_j^\top p^{(1)}$. Marginalization gives
\begin{equation}\label{eq:marginal-model}
z_j\mid\mu,\lambda,\sigma^2
\sim N\left(b_j\mu,\sigma^2c_j(\lambda)\right),\quad
c_j(\lambda)=1+\frac{1}{\lambda\delta_j}.
\end{equation}
This is a Gaussian residual/marginal-likelihood interpretation,
not a new form of PLS. It is related to likelihood-based
Whittaker--Henderson smoothing \cite{biessy2026wh}.
Profiling the variance $\sigma^2$ for a fixed numerical $\mu$
leaves, up to an additive constant,
\begin{equation}\label{eq:likelihood-profile}
J(\lambda;\mu)=r\log\left(
\frac{1}{r}\sum_{j=1}^{r}
\frac{(z_j-b_j\mu)^2}{c_j(\lambda)}
\right)+\sum_{j=1}^{r}\log c_j(\lambda).
\end{equation}
An independently developed branch of the package,
\texttt{feature/mle-smoothness-and-data-mu-20261021},
adds this profiled Gaussian score for fixed $\mu$ or
the conditional plug-in $\mu=\overline{D_dy}$.
That plug-in must not be mislabeled as a joint MLE.

\subsection{A closed-form estimator of $\mu$ conditional on $\lambda$}
Let $w_j(\lambda)=c_j(\lambda)^{-1}
=\lambda\delta_j/(1+\lambda\delta_j)$.
For a fixed positive $\lambda$, the generalized
least-squares optimum satisfies
\begin{equation}\label{eq:mu-profile-mle}
\widehat\mu(\lambda)
=\frac{\sum_{j=1}^{r}w_jb_jz_j}
{\sum_{j=1}^{r}w_jb_j^2}.
\end{equation}
The denominator is positive because $D_dp^{(1)}=\mathbf1$.
Equation~\eqref{eq:mu-profile-mle} is also obtained by
minimizing the PLS objective jointly in $\tau$ and $\mu$
for a fixed $\lambda$. Substitution in
\eqref{eq:likelihood-profile} gives the one-dimensional
concentrated score
$J(\lambda;\widehat\mu(\lambda))$,
whose optimization produces a joint profile-likelihood
candidate $(\widehat\lambda,\widehat\mu)$.
Such a solution depends on the Gaussian variance model,
may lie near a boundary, and must be checked for
identifiability and other local minima. This paper provides
the equation and a separately reproducible numerical test;
joint $\mu$ profiling is not yet an exposed release API.

In fact, taking the limit $\lambda\downarrow0$
using $w_j\sim\lambda\delta_j$ proves
\begin{equation}\label{eq:mu-guerrero-link}
\lim_{\lambda\downarrow0}\widehat\mu(\lambda)
=\frac{(p^{(1)})^\top Q_d y}{(p^{(1)})^\top Q_dp^{(1)}}
=\frac{\mathbf1^\top D_dy}{n-d}
=\overline{D_dy}.
\end{equation}
For positive $\lambda$ the two expressions differ in general.
Estimating $\mu$ adds a rank-one adjustment to the
data-to-trend linear operator, so the fixed-$\mu$ effective
degrees of freedom and ordinary leave-one-out denominators
are not automatically valid after joint estimation.
Prior literature already considers nonzero difference
means \cite{guerrero2007penalized,cortes2017smoothness};
the implementation contribution here is an explicit,
testable profiling workflow, not a claim to have first
introduced $\mu$ or likelihood selection.
